\chapter*{完整源码清单}
\addcontentsline{toc}{chapter}{完整源码清单}

本附录不是节选代码，而是从 \filepath{books/compute-systems-engine-code/labs} 直接读入的完整工程源码。读者在 PDF 中看到的清单，应当和本地可构建工程保持同源；如果源码修改，重建书籍后清单会随之更新。

源码阅读顺序不要从最长的实现文件开始。先读 CMake 和 README，确认工程边界；再读公共头文件，确认类型和函数合同；然后读实现和 probe，观察机制如何跑起来；最后读测试和 benchmark，检查哪些行为已经被固定，哪些只是性能 smoke。

\section*{一、构建入口与 Reference Pipeline}
\addcontentsline{toc}{section}{一、构建入口与 Reference Pipeline}

这一组文件给出代码实践卷的顶层构建入口，以及第一条主线 \filepath{reference_pipeline} 的完整源码。它从输入身份、解析诊断、reference 结果、report、manifest、split 和 merge 开始，训练“先定义正确结果，再谈优化”的工程习惯。逐文件解释见第三章；本附录只负责保证源码完整同源。

\lstinputlisting[language=bash,caption={代码实践卷 README.md}]{../../README.md}

\lstinputlisting[language=bash,caption={代码实践卷 Makefile}]{../../Makefile}

\lstinputlisting[language=bash,caption={代码实践卷顶层 CMakeLists.txt}]{../../CMakeLists.txt}

\lstinputlisting[language=bash,caption={代码实践卷 LaTeX Makefile}]{Makefile}

\lstinputlisting[language=bash,caption={reference pipeline README.md}]{../../labs/reference_pipeline/README.md}

\lstinputlisting[language=bash,caption={reference pipeline CMakeLists.txt}]{../../labs/reference_pipeline/CMakeLists.txt}

\lstinputlisting[language=C++,caption={include/cse/reference\_pipeline.hpp}]{../../labs/reference_pipeline/include/cse/reference_pipeline.hpp}

\lstinputlisting[language=C++,caption={src/reference\_pipeline.cpp}]{../../labs/reference_pipeline/src/reference_pipeline.cpp}

\lstinputlisting[language=C++,caption={reference pipeline src/main.cpp}]{../../labs/reference_pipeline/src/main.cpp}

\section*{二、Compute Systems 合同层}
\addcontentsline{toc}{section}{二、Compute Systems 合同层}

这一组文件给出第二条主线 \filepath{compute_systems} 的构建入口、阅读地图和公共头文件。公共头文件是合同层：每个 struct 都把一个章节机制压成可观察字段，每个函数签名都说明输入规模、边界和返回结果。

\lstinputlisting[language=bash,caption={compute systems README.md}]{../../labs/compute_systems/README.md}

\lstinputlisting[language=bash,caption={compute systems CMakeLists.txt}]{../../labs/compute_systems/CMakeLists.txt}

\lstinputlisting[language=C++,caption={include/compsys/chapter\_models.hpp}]{../../labs/compute_systems/include/compsys/chapter_models.hpp}

\lstinputlisting[language=C++,caption={include/compsys/parallel\_reduce.hpp}]{../../labs/compute_systems/include/compsys/parallel_reduce.hpp}

\lstinputlisting[language=C++,caption={include/compsys/sync\_primitives.hpp}]{../../labs/compute_systems/include/compsys/sync_primitives.hpp}

\lstinputlisting[language=C++,caption={include/compsys/wait\_channels.hpp}]{../../labs/compute_systems/include/compsys/wait_channels.hpp}

\lstinputlisting[language=C++,caption={include/compsys/task\_runtime.hpp}]{../../labs/compute_systems/include/compsys/task_runtime.hpp}

\lstinputlisting[language=C++,caption={include/compsys/futex\_lab.hpp}]{../../labs/compute_systems/include/compsys/futex_lab.hpp}

\section*{三、Compute Systems 实现和探针}
\addcontentsline{toc}{section}{三、Compute Systems 实现和探针}

这一组文件是第二条主线的完整实现与命令行探针。实现文件负责让章节模型、并行 reduce、同步原语、wait channel、任务运行时和 futex 实验真正运行；probe 文件负责打印稳定输出，让书中结论、命令行输出和测试断言可以互相对照。

\lstinputlisting[language=C++,caption={src/chapter\_models.cpp}]{../../labs/compute_systems/src/chapter_models.cpp}

\lstinputlisting[language=C++,caption={src/parallel\_reduce.cpp}]{../../labs/compute_systems/src/parallel_reduce.cpp}

\lstinputlisting[language=C++,caption={src/sync\_primitives.cpp}]{../../labs/compute_systems/src/sync_primitives.cpp}

\lstinputlisting[language=C++,caption={src/wait\_channels.cpp}]{../../labs/compute_systems/src/wait_channels.cpp}

\lstinputlisting[language=C++,caption={src/task\_runtime.cpp}]{../../labs/compute_systems/src/task_runtime.cpp}

\lstinputlisting[language=C++,caption={src/futex\_lab.cpp}]{../../labs/compute_systems/src/futex_lab.cpp}

\lstinputlisting[language=C++,caption={compute systems src/main.cpp}]{../../labs/compute_systems/src/main.cpp}

\lstinputlisting[language=C++,caption={src/chapter\_models\_probe.cpp}]{../../labs/compute_systems/src/chapter_models_probe.cpp}

\lstinputlisting[language=C++,caption={src/wait\_channels\_probe.cpp}]{../../labs/compute_systems/src/wait_channels_probe.cpp}

\lstinputlisting[language=C++,caption={src/task\_runtime\_probe.cpp}]{../../labs/compute_systems/src/task_runtime_probe.cpp}

\lstinputlisting[language=C++,caption={src/futex\_lab\_probe.cpp}]{../../labs/compute_systems/src/futex_lab_probe.cpp}

\section*{四、测试和 Benchmark}
\addcontentsline{toc}{section}{四、测试和 Benchmark}

这一组文件是进入正文结论的门槛。测试固定 reference pipeline 和 compute systems 的语义边界；benchmark 只提供性能入口，不替代正确性测试。新增优化前，先跑测试，再跑 benchmark smoke，最后才写性能报告。

\lstinputlisting[language=C++,caption={reference pipeline tests/reference\_pipeline\_tests.cpp}]{../../labs/reference_pipeline/tests/reference_pipeline_tests.cpp}

\lstinputlisting[language=C++,caption={bench/reference\_pipeline\_bench.cpp}]{../../labs/reference_pipeline/bench/reference_pipeline_bench.cpp}

\lstinputlisting[language=C++,caption={compute systems tests/compsys\_tests.cpp}]{../../labs/compute_systems/tests/compsys_tests.cpp}
